\documentclass[11pt,reqno]{amsart}
\usepackage{amsmath,amssymb,graphicx,amscd}
\usepackage{mathrsfs}
\usepackage{enumerate}
% pour les S des spheres
\usepackage{mathrsfs}
\usepackage{dsfont}
\usepackage{stmaryrd}
\usepackage{wrapfig}
\usepackage{pstricks}

% ----------------------------------------------------------------
\vfuzz2pt % Don't report over-full v-boxes if over-edge is small
\hfuzz2pt % Don't report over-full h-boxes if over-edge is small

\newcommand{\prob}{{\mathbb{P}}}

\renewcommand{\S}{\mathscr{S}}  % complex unit sphere

\newcommand{\CC}{{\mathbb C}}
\newcommand{\HH}{{\mathbb H}}
\newcommand{\RR}{{\mathbb R}}
\newcommand{\QQ}{{\mathbb{Q}}}
\newcommand{\NN}{{\mathbb{N}}}
\newcommand{\ZZ}{{\mathbb{Z}}}
\renewcommand{\P}{{\mathbb{P}}}
\newcommand{\sgn}{\operatorname{{sgn}}}
 \DeclareMathOperator{\Prob}{\mathbb{P}}   %probability
 \DeclareMathOperator{\E}{\mathbb{E}}      %expectation
 \DeclareMathOperator{\cov}{cov}             %logarithmic integral
\newcommand{\1}{\mathds{1}}

 \newcommand{\I}{1\!\!1}                   %Indicator function


 \newcommand{\rd}{{\mathrm{d}}}            %differential
 \newcommand{\ii}{{\mathrm{i}}}
  \newcommand{\Id}{{\mathrm{Id}}}
\newcommand{\law}{\overset{\mbox{\rm \scriptsize law}}{=}}
\newcommand{\convlaw}{\overset{\mbox{\rm \scriptsize law}}{\longrightarrow}}


\renewcommand{\Re}{{\mathfrak{Re}}}
\renewcommand{\Im}{{\mathfrak{Im}}}

% THEOREMS -------------------------------------------------------

\newtheorem{thm}{Theorem}[section]
\newtheorem{cor}[thm]{Corollary}
\newtheorem{lem}[thm]{Lemma}
\newtheorem{prop}[thm]{Proposition}
\theoremstyle{definition}
\newtheorem{defn}[thm]{Definition}
\theoremstyle{remark}
\newtheorem*{rem}{Remark}



\numberwithin{equation}{section}


\begin{document}

~\vspace{-1cm}



\begin{center}
{\bf Probability, homework 9, due November 18.}
\end{center}
\vspace{0.5cm}
\small







\noindent{\bf Exercise 1.} For fixed $p,q\in[0,1]$,  consider a Markov chain $X$ with two states $\{1,2\}$, with transition matrix
$$
\pi=(\pi(i,j))_{1\leq i,j\leq 2}=\left(
\begin{array}{cc}
1-p&p\\
q&1-q
\end{array}
\right)
$$
\begin{enumerate}[(i)]
\item For which $p,q$ is the chain irreducible? Aperiodic?
\item What are the invariant probability measures of $X$?
\item Compute $\pi^{(n)}$, $n\geq 1$.
\item When $X$ is irreducible,  for this invariant probability measure $\mu$, calculate
\begin{align*}
&d_1(n):=\frac{1}{2}\left(|\P_1(X_n=1)-\mu(1)|+|\P_1(X_n=2)-\mu(2)|\right)\\
&d_2(n):=\frac{1}{2}\left(|\P_2(X_n=1)-\mu(1)|+|\P_2(X_n=2)-\mu(2)|\right)
\end{align*}
where $\P_x$ means the chain starts at $x$.
\end{enumerate}



\noindent{\bf Exercise 2.} Consider a Markov chain $X$ with state space $\mathbb{N}$ and transition matrix 
$$
\pi(0,0)=r_0,\ \pi(0,1)=p_0,  {\rm and}\ \forall i\geq 1,\ \pi(i,i-1)=q_i,\ \pi(i,i)=r_i,\ \pi(i,i+1)=p_i,
$$
with $p_0,r_0>0$, $p_0+r_0=1$ and for all $i\geq 1$, $p_i,q_i>0$, $p_i+q_i+r_i=1$. 
Prove that the chain is irreducible, aperiodic.  Give a necessary and sufficient condition for the chain to have an invariant probability measure.\\















\noindent{\bf Exercise 3.} Consider a Markov chain $X$ with state space $\{0,1,\dots,n\}$ and transition matrix 
\begin{align*}
\pi(0,k)&=\frac{1}{2^{k+1}},\ 0\leq k\leq n-1,\ \pi(0,n)=\frac{1}{2^n}\\
\pi(k,k-1)&=1,\ 1\leq k\leq n-1,\ \pi(n,n)=\pi(n,n-1)=\frac{1}{2}.
\end{align*}
\begin{enumerate}[(i)]
\item Prove that the chain has a unique invariant probability measure $\mu$ and calculate it.
\item Prove that for any $0\leq x_0\leq n-1$,  $\pi^{(x_0+1)}(x_0,\cdot)=\mu$.
\item Prove that for any $0\leq x_0\leq n$,  $\pi^{(n)}(x_0,\cdot)=\mu$.
\item For any $t\geq 1$, calculate 
$$
d(t):=\frac{1}{2}\sum_{x=0}^n\left|
\pi^{(t)}(n,x)-\mu(x)
\right|,
$$ 
and plot $t\mapsto d(t)$.
\end{enumerate}




\noindent{\bf Exercise 4.} Let $(G,\cdot)$ be a group,  $\mu$ a probability measure on $G$ and X the Markov chain such that $\pi(g,h\cdot g)=\mu(h)$. We call such a process $X$ a random walk on $G$ with jump kernel $\mu$.
\begin{enumerate}[(i)]
\item Explain why the usual random walk on $\mathbb{Z}^d$ is such process. Same question for the usual random walk on $(\mathbb{Z}/n\mathbb{Z})^d$, $n\geq 1$.
\item Consider the following shuffling of a deck of $n\geq 2$ cards: pick two such distinct cards uniformly at random and exhange their positions in the deck. Show that this is also an example of a random walk on a group.
\item Let $\mathcal{H}=\{h_1\cdot h_2\cdot\dots\cdot h_n,\mu(h_i)>0, 1\leq i\leq n,n\in\mathbb{N}\}$.  Discuss irreductibility of $X$ depending on $\mathcal{H}$.
\item If $X$ is irreducible on finite $G$, what are the invariant probability measures? What if $G$ is not finite?
\item Make some search to define a reversible Markov chain. In the context of this exercise, show that $X$ is reversible if and only if $\mu(h)=\mu(h^{-1})$ for any $h\in G$.
\item Give an example of an irreducible random walk on a group which is not reversible.
\end{enumerate}


\noindent{\bf Exercise 5.}
An ant walks on a round clock, starting at 0, up to the moment it has visited all numbers.  At each second,  it walks either clockwise or counterclockwise,   with probability $1/2$ to a neighbouring number,  and through independent steps.  Let $X$ be the final position of the ant. Prove it is equidistributed on $\{1,2,\dots,11\}$.\\
































\end{document} 